% P11 paper module: R35 contraction no-go for R_{T0}, exact resolvent remainder,
% and unconditional rescaled Neumann series repair of R34.
% Records the two-window/two-prime projection argument showing \|R_{T_0}\|>1 for
% every T_0>log 3, replaces the illegitimate divergent-tail language of R34 by an
% exact finite resolvent remainder identity valid for every bounded A>=0, and gives
% an unconditional rescaled Neumann series for B_{T_0} that does not require any
% spectral gap. No polar-gauge, terminal-transport, Object-X, or RH consequence is
% drawn; R30-F and the R32-F(ii) fingerprint route remain open.

\subsection{Contraction no-go for \(R_{T_0}\) and repair of the R34 Neumann tail}
\label{subsec:o3ag-contraction-nogo}

R34 (module chain above) attempted a Neumann-series route to the regularity of
\(g_{R,S}=-B_{T_0}H_{T_0}E_{S,T_0}j_{R,S}\) via
\(B_{T_0}=(I+R_{T_0}^*R_{T_0})^{-1}\), conditional on a spectral gap
\(\|R_{T_0}\|<1\) (Open Problem R34-C).  This subsection closes that gap question in
its large-horizon contraction form, negatively, and repairs an imprecise
divergent-series formulation in the R34 audit.  It records neither Object~X nor
polar-gauge nor terminal-transport nor RH consequences.

\subsubsection*{Two-prime contraction no-go}

For \(p\in\{2,3\}\) put \(\tau_p:=\tfrac12\log p\).

\begin{theorem}[Large-horizon contraction no-go]
\label{thm:o3ag-contraction-nogo}
Fix \(T_0>\log3\) and choose
\[
 0<\varepsilon<\min\Bigl\{T_0-\log3,\ \tfrac{\tau_2}{3},\ \tfrac{\tau_3-\tau_2}{3}\Bigr\}.
\]
Let \(f\in L^2(-T_0,T_0)\), \(\|f\|=1\), \(\operatorname{supp}f\subset(-\varepsilon,\varepsilon)\),
and for \(p=2,3\) put \(W_{p,\pm}:=(\pm\tau_p-\varepsilon,\pm\tau_p+\varepsilon)\).  Let
\(\Pi_p\) denote the orthogonal projection of the target field of \(R_{T_0}\) onto the
\(p\)-prime component intersected with \(W_{p,+}\cup W_{p,-}\).  Then
\[
 \boxed{
 \|\Pi_pR_{T_0}f\|^2=2(\log p)\,p^{-1/2}\Bigl(1-\frac1p\Bigr),\qquad p\in\{2,3\},
 }
\]
and, since the \(p=2\) and \(p=3\) target sectors are orthogonal,
\[
 \boxed{
 \|R_{T_0}f\|^2\ge\|\Pi_2R_{T_0}f\|^2+\|\Pi_3R_{T_0}f\|^2
 =\frac{\log2}{\sqrt2}+\frac{4\log3}{3\sqrt3}\approx1.335830.
 }
\]
Hence
\[
 \boxed{
 \|R_{T_0}\|\ge\sqrt{\frac{\log2}{\sqrt2}+\frac{4\log3}{3\sqrt3}}\approx1.15578>1
 \qquad\text{for every }T_0>\log3.
 }
\]
\end{theorem}

\begin{proof}
By~\eqref{eq:eta-pk} at \(k=1\), \(\eta_{p,1}=\sqrt{p-1}\,p^{-1/2}\psi_{p,0}\), so
\(\|\eta_{p,1}\|^2=(p-1)/p=1-1/p\).  For \(|u|<\tau_p+\varepsilon\le\log3+\varepsilon<T_0\)
(using \(\varepsilon<T_0-\log3\)), the depth~\eqref{eq:martingale-depth} satisfies
\(J_{p,T_0}(u)\ge1\), so by~\eqref{eq:source-cutoff}--\eqref{eq:projected-mark} the full
\(k=1\) mark survives on both windows \(W_{p,\pm}\), i.e.\ \(\mathsf Q_{T_0}(u)\eta_{p,1}=\eta_{p,1}\)
there.

The rest operator's \(k=1\) summand for prime \(p\) in~\eqref{eq:rest-operator} is
\(\sqrt{\log p}\,p^{-1/4}(D_{\log p}E_{T_0}f)(u)\otimes\mathsf Q_{T_0}(u)\eta_{p,1}\), and
\((D_{\log p}f)(u)=f(u-\tau_p)-f(u+\tau_p)\).  Since
\(\operatorname{supp}f\subset(-\varepsilon,\varepsilon)\) and \(\varepsilon<\tau_p/3\), the
two shifted copies \(f(\cdot-\tau_p)\) (supported in \(W_{p,+}\)) and \(f(\cdot+\tau_p)\)
(supported in \(W_{p,-}\)) have disjoint supports, so
\[
 \|\Pi_p(D_{\log p}E_{T_0}f)\|^2=\|f(\cdot-\tau_p)\|^2+\|f(\cdot+\tau_p)\|^2=2.
\]
No \(k\ge2\) shift of the same prime sector reaches \(W_{p,\pm}\): for \(p=2\), the
\(k=2\) shift is centered at \(2\tau_2\), and \(2\tau_2-\varepsilon>\tau_2+\varepsilon\)
since \(\varepsilon<\tau_2/3\); for \(p=3\), analogously with \(\varepsilon<\tau_3/3\).
Combined with orthogonality of different prime sectors, multiplying by the coefficient
\((\log p)p^{-1/2}\) and the mark norm \(1-1/p\) gives the boxed identity for
\(\|\Pi_pR_{T_0}f\|^2\); numerically \(p=2\): \((\log2)/\sqrt2\approx0.49013\), \(p=3\):
\(4(\log3)/(3\sqrt3)\approx0.84571\).  Orthogonality of the \(p=2\) and \(p=3\) target
sectors and Pythagoras give the stated lower bound on \(\|R_{T_0}f\|^2\); since
\(\|f\|=1\), taking square roots gives the norm bound.
\end{proof}

\begin{corollary}[The contraction route is dead for large horizons]
\label{cor:o3ag-nogo}
For every \(T_0>\log3\), hypothesis~\eqref{eq:r34-gap} of Proposition R34-B fails:
\(\|R_{T_0}\|<1\) does not hold.  In particular the full unscaled Neumann series
\(\sum_{n\ge0}(-1)^n(R_{T_0}^*R_{T_0})^nw\) diverges in operator norm on this range and
cannot be used, even termwise, to transfer regularity from \(w\) to \(B_{T_0}w\).
\end{corollary}

\begin{remark}[Scope]
Theorem~\ref{thm:o3ag-contraction-nogo} does not decide \(\|R_{T_0}\|<1\) for
\(0<T_0\le\log3\).  It says nothing about boundedness or positivity of \(B_{T_0}\),
which hold unconditionally for every \(T_0\) (Lemma~\ref{lem:o3ag-rescaled-neumann}
below).  It rules out only this specific contraction route, not other possible paths to
regularity of \(g_{R,S}\).
\end{remark}

\subsubsection*{Exact resolvent remainder and unconditional rescaled series}

\begin{lemma}[Exact finite resolvent remainder]
\label{lem:o3ag-exact-remainder}
For every bounded \(A\ge0\), every \(N\ge0\), and \(B:=(I+A)^{-1}\),
\[
 \boxed{
 B=\sum_{n=0}^N(-A)^n+(-1)^{N+1}A^{N+1}B,
 }
\]
with no norm hypothesis on \(A\).
\end{lemma}

\begin{proof}
Induction on \(N\).  For \(N=0\): \((I+A)B=I\) gives \(B=I-AB\).  Assuming the identity
at \(N\), substitute \(B=I-AB\) into the remainder term to advance to \(N+1\).
\end{proof}

This is the legitimate replacement for the informal ``tail''
\(\sum_{n>N}(-1)^n(R_{T_0}^*R_{T_0})^nw\) used in the R34 audit; at
\(\|A\|\ge1\) that infinite sum need not converge and must not be invoked, whereas the
finite remainder \((-1)^{N+1}A^{N+1}B\) is exact and unconditional for every \(N\).

\begin{lemma}[Unconditional rescaled Neumann series]
\label{lem:o3ag-rescaled-neumann}
For every bounded \(A\ge0\) and every \(M\ge\|A\|\), put
\(Q_M:=(MI-A)/(1+M)\).  Then
\[
 0\le Q_M\le\frac{M}{1+M}I,
 \qquad
 \|Q_M\|<1,
\]
and
\[
 \boxed{
 B=(I+A)^{-1}=\frac1{1+M}\sum_{n=0}^\infty Q_M^n
 }
\]
in operator norm.
\end{lemma}

\begin{proof}
Since \(0\le A\le MI\), \(0\le MI-A\le MI\), giving the stated bound on \(Q_M\) and hence
\(\|Q_M\|<1\), so \(\sum_nQ_M^n=(I-Q_M)^{-1}\) converges in norm.  Direct computation
gives \(I-Q_M=(I+A)/(1+M)\), so \((I-Q_M)^{-1}=(1+M)B\), i.e.\
\(B=(1+M)^{-1}\sum_nQ_M^n\).
\end{proof}

\begin{remark}[Spectral gap is not a necessary gate]
Lemma~\ref{lem:o3ag-rescaled-neumann} shows that a convergent series representation of
\(B_{T_0}\) exists for every bounded \(A=R_{T_0}^*R_{T_0}\ge0\), independent of whether
\(\|R_{T_0}\|<1\).  Open Problem R34-C, read as ``does a Neumann-type representation of
\(B_{T_0}\) exist,'' is therefore answered trivially and unconditionally.  It must be
reread as asking whether the iterates \(A^{N+1}\) (Lemma~\ref{lem:o3ag-exact-remainder})
or \(Q_M^n\) (Lemma~\ref{lem:o3ag-rescaled-neumann}) carry regularity or support
structure usable for R32-F(i); this is a smoothing/fingerprint question about the
iterates, not a spectral-gap question, and is not decided here.
\end{remark}

\begin{remark}[Split status of R34-C]
\label{rem:o3ag-r34c-split}
R34-C is accordingly split into: (a) the naive unscaled contraction route
\(\|R_{T_0}\|<1\) --- closed negatively by Theorem~\ref{thm:o3ag-contraction-nogo} for
every \(T_0>\log3\); (b) existence of \emph{some} convergent series for \(B_{T_0}\) ---
answered affirmatively and unconditionally by Lemma~\ref{lem:o3ag-rescaled-neumann};
(c) regularity/fingerprint structure of the iterated remainder --- genuinely open, and
unaffected by (a) or (b).  Only (c) survives as live mathematical content.
\end{remark}

\begin{remark}[Consequence for the route map]
\label{rem:o3ag-routemap}
Because (a) is dead for \(T_0>\log3\), Strategy (i) of Open Problem R32-F is not viable
in its spectral-gap form for the large terminal horizons of ultimate interest.  This
motivates opening R32-F Strategy (ii) --- the localized annular range/annihilator
fingerprint of \(\Sigma_{T_0}=H_{T_0}B_{T_0}H_{T_0}^*\) --- as the next audit target.
That route must first decide whether
\[
 \ker(H_{T_0}E_{\mathcal A})=\{0\}?
\]
for the annulus zero-extension \(E_{\mathcal A}\) before attempting to construct any
annihilator witness; existence of a localized annihilator is not assumed here and
remains to be proved or refuted.  It does not otherwise change the mathematical status
of R30, R30-F, R31, R32, or R32-F.  In particular R30-F remains \(?[O]\).  No
polar-gauge, terminal-transport, Object-X, or RH statement is made in this subsection.
\end{remark}
